% ЕМТ 2019, VIII клас — проверен LaTeX препис на липсващите/непълните задачи 8.1 и 8.3 от официалните файлове в Klasirane.com.
% Условия: https://klasirane.com/api/competitions/EMT/All/2019/8%20%D0%BA%D0%BB/probs
% SHA-256 (условия): bfc34dc349672e81d092cfc2e8f8c050271342a991e87d2dd3297323c869aea0
% Решения: https://klasirane.com/api/competitions/EMT/All/2019/8%20%D0%BA%D0%BB/sol
% SHA-256 (решения): 6395ca5590260fb23470f8ad3d8e1e7af19f8b067c9b52814855c087f09341ec
% Mathpix Snip (условия): 2fb98f33-33c6-4e6d-985d-d861135f7f26
% Mathpix Snip (решения): 55a37411-eaa3-4c42-a913-c5a7f102de0f
% Формулировките, решенията и точковите оценки са сверени визуално с официалните PDF страници; задачите нямат фигури.
% Изброените случаи в оценяването на 8.3 са запазени дословно от източниковия PDF.

Задача 8.1. Да се намерят всички цели стойности на $x$, за които стойността на израза $M=|x^3-2x^2-10x+8|$ е просто число.

Решение. Разлагаме $M=|x-4|\left|x^{2}+2x-2\right|$. $M$ ще е просто число, ако единият множител е равен на 1, а другият е просто число.

1 сл. $|x-4|=1 \Rightarrow x \in\{3,5\}$. При $x=3$ получаваме $M=13$, което е просто. При $x=5$ получаваме $M=33$, което не е просто.

2 сл. $\left|x^{2}+2x-2\right|=1 \Rightarrow x^{2}+2x-2=\pm1$. От първото уравнение намираме $x=1$ или $x=-3$ и, съответно, $M=3$ или $M=7$, които са прости числа. Второто уравнение е еквивалентно на $(x+1)^{2}=2$, което няма решение в цели числа.

Окончателно, решенията са $x=1$ и $x=\pm3$.

Оценяване. (6 точки) 2 т. – за разлагането на $M$; по 2 т. – за решаването на всеки от двата случая.

Задача 8.3. Да се намерят всички естествени числа $p$ и $n$, такива че $p$ е просто и $p^3+5n-1=n(p^2+2p+6n)$.

Решение. Условието е равносилно с $p\left(p^{2}-pn-2n\right)=(2n-1)(3n-1)$, така че $p$ дели $2n-1$ или $3n-1$. Ако $2n-1\geq2p$ или $3n-1\geq3p$, то лявата страна е отрицателна, а дясната – положителна. Остава да проверим случаите:
\begin{itemize}
\item $p=2n-1$, което не води до естествени решения;
\item $p=3n-1$, което не води до естествени решения;
\item $2p=3n-1$, което води до
$$
\begin{gathered}
p^{2}-pn-2n=4n-2,\\
9n^{2}-6n+1-2n(3n-1)=24n-8,\\
3n^{2}-28n+9=0,\\
(3n-1)(n-9)=0,
\end{gathered}
$$
чието естествено решение е само $n=9$. Оттук $p=13$ и условията са изпълнени.
\end{itemize}

Оценяване. (7 точки) 2 т. – за доказателство, че $p$ дели $2n-1$ или $3n-1$; по 1 т. – за пълно изследване на всеки от случаите $2n-1 \geq 2p$, $2n-1=p$, $3n-1 \geq 2p$, $3n-1=p$.
